\chapter{Conclusion}
\label{c:conclusion}

This thesis examined the state-of-the-art in intentionally vulnerable Android app design as found in the OWASP MASTG reference apps, and conducted a literature review to understand how the overarching OWASP MAS project is used in the academic world. Building on this knowledge, a set of intentionally vulnerable Android apps was developed to expand the MASTG reference app catalog and improve its function as a teaching resource and as a benchmark dataset for different types of mobile app security testing tools. 

The literature review showed that the main use-case of the OWASP MAS project in academic works consists of using the MASTG for evaluating mobile apps based on the quality of their security and whether they meet the requirements laid out by the MASVS. Especially the MASTG penetration testing guide was used as a tool for evaluating app security by numerous authors. Multiple papers used the MASVS vulnerabilities categories to map and classify real-world Android security vulnerabilities. The MASTG reference apps were used for testing and evaluating newly developed tools dealing with vulnerability detection. Overall, the literature review shows that the academic world uses the OWASP MAS project mainly as a framework and tool for evaluating mobile app security.

The analysis of the existing OWASP MASTG reference apps showed that only a few have documentation on the vulnerabilities they include. Furthermore, only a small portion of the vulnerabilities as described by the OWASP MASWE find any representation in the reference apps. Most vulnerability types have no reference implementation in the apps, whereas a small group of them is implemented multiple times in different ways. The ten most frequently implemented vulnerability types make up almost half of all implemented vulnerabilities, despite there being 117 different vulnerability types described by the OWASP MASWE. The same trend can be found when examining the vulnerability categories themselves. A few categories, such as MASVS-PLATFORM or MASVS-STORAGE, have most of their respective vulnerability types implemented at least once in a MASTG reference app, while other categories see very little representation in the reference app catalog. Not one of the nine vulnerabilities described by MASVS-PRIVACY is included in any of the reference apps.

The apps we developed are designed to reflect the natural ordering of vulnerabilities done by the MASWE. Each MASVS category makes up its own app, which are in turn a part of the whole project in the form of modules. For each app, the vulnerabilities that comprise its designated MASVS category are grouped into separate packages, with one package per implemented MASWE vulnerability. Further extending this architecture with a shared library module containing common attributes, and extensive use of custom templates and inheritance, allows for a design which can easily be extended with new vulnerabilities and categories while producing minimal scaling costs. The vulnerability selection was focused on MASWE vulnerability types that have not been implemented in any of the existing MASTG reference apps. Instead of working on multiple MASVS categories at once, the vulnerability categories were implemented one by one.

The three apps developed are called \texttt{maswe\_storage}, \texttt{maswe\_crypto} and \texttt{maswe\_platform}. Together they contain 28 vulnerability types, 20 of which have not been implemented in any of the existing MASTG reference apps. All 6 and all 18 vulnerability types described by MASVS-STORAGE and MASVS-CRYPTO have been implemented. Through the addition of our apps, the total coverage of MASWE vulnerability types has been increased from 32.5\% to 49.6\%. All implemented vulnerabilities are mapped to their MASWE counterparts and have extensive documentation on where they can be found, how they can be exploited, and sometimes how they can be fixed. The apps developed are an important milestone in building a unified, well structured and thoroughly documented MASTG reference apps benchmark.

A recurring challenge during the development of the apps was the limited documentation by OWASP on the MASWE vulnerabilities. The OWASP MAS project is constantly evolving, as new vulnerabilities become more relevant and older ones become obsolete. This leads to many of the vulnerabilities described by MASWE still being in BETA and having very little documentation on how they should be implemented and what they should include. This complicates the development process, since the vulnerabilities themselves need to have adequate depth for them to be useful as an educational resource and security tool testing dataset.

This thesis can be expanded upon in many different ways. Half of the vulnerability types described by MASWE still have no representation in the MASTG reference apps. Now that the apps are developed, extensive benchmarking of both static and dynamic security tools can be carried out in order to better understand the mobile app security testing tools landscape. A complimentary set of iOS MASTG reference apps can be developed, making use of the architecture and design built to branch out to the iOS side of the OWASP MAS project.

This work greatly expands the coverage of MASWE vulnerability types made by the existing MASTG reference applications catalog. Through this, it contributes to making Android mobile security more accessible and easy to learn. The extensive documentation on the complete set of implemented MASVS-STORAGE and MASVS-CRYPTO vulnerabilities builds towards a more systematic environment for mobile security research and education. The applications developed serve as the foundation for a complete OWASP MAS benchmark for evaluating mobile application security analysis tools. The scalable and modular architecture on which the apps have been developed enables them to change alongside the Android mobile security landscape and the OWASP MAS project as they continue to evolve, which allows for easy maintenance and continuous growth as a completely open-source benchmark for mobile application security.